\documentclass[../dayan_kernel_manual.tex]{subfiles}

\begin{document}

\section{RISC-V 与对象 ABI}\label{sec:kernel-080}

\subsection{数据与调用约定}

Device ABI 为 RV64 \code{lp64}。\code{char}/\code{short}/\code{int} 分别为
8/16/32 bit；\code{long}、指针、\code{size_t} 和 \code{uintptr_t} 为 64 bit。公开
函数边界和上下文帧保持 16-byte 栈对齐。

\code{lp64} 使用整数参数寄存器。ISA 含 F 扩展，函数内部可以执行单精度指令，但函数
参数和返回值仍按整数 ABI 传递。dycc wrapper 把 \code{float}/\code{double} 位模式
放入 64-bit 参数槽；直接使用低层 launch API 时由调用者完成同样的编码。

\begin{dayantable}{保留寄存器}{tab:kernel-080}{X[18,l] X[22,l] X[60,l]}
  寄存器 & 名称 & 用途 \\
  \code{x2/sp} & stack pointer & 每 proc 独立，边界保持 16-byte 对齐。 \\
  \code{x3/gp} & \code{shared} & 当前 proc 的 Host coherent allocation 基址。 \\
  \code{x4/tp} & \code{current} & 当前 \code{struct thread}；不能用于 ELF TLS。 \\
  \code{x10--x17} & \code{a0--a7} & 新 proc 的 0--8 个机器字参数。 \\
\end{dayantable}

固件不定义 \code{__global_pointer\$}，\code{gp} 不作为 small-data pointer。固件和
算子均使用 \code{-mno-relax} 与 \code{-msmall-data-limit=0}，不得产生 GP-relative、
TP-relative 或 TLS relocation。

\subsection{proc 与切换帧}

新 proc 使用 96-byte 启动帧：

\begin{dayantable}{proc 启动帧}{tab:kernel-081}{X[20,l] X[22,l] X[58,l]}
  offset & 字段 & 含义 \\
  0 & \code{entry} & kernel entry。 \\
  8 & \code{ra} & \code{proc_exit}。 \\
  16--72 & \code{a0--a7} & 8 个参数槽。 \\
  80 & \code{shared} & 写入 \code{gp} 的基址。 \\
  88 & reserved & 0。 \\
\end{dayantable}

bootstrap 入口为 \code{proc_bootstrap + (8 - argc) * 4}，只执行所需的参数加载指令。
上下文切换帧为 128 bytes，依次保存 \code{ra}、\code{s0--s11}、\code{gp}、
\code{fcsr} 和 reserved word。

\code{struct thread} 为 56 bytes、8-byte 对齐。type/state/id 位于 offset 0/1/2，
parent 和两字节 intrusive node 位于 4--6，hardware-token outstanding 标志位于 7，
sp/pc 位于 8/16，Context1/Context2 位于 24/32，token extension 位于 36，16-byte 名称
从 37 开始。派生对象把 thread 放在 offset 0。

\subsection{公开结构体}

\begin{dayantable}{Device 公开结构体布局}{tab:kernel-082}{X[31,l] X[17,l] X[52,l]}
  类型 & 大小/对齐 & 字段 offset \\
  \code{dayan_resource_t} & 1/1 & 按值传递的 8-bit 逻辑编号；不含状态或指针。 \\
  \code{dayan_resource_mask_t} & 8/8 & bit 0--31 为 Local，bit 32--63 为 Shared。 \\
  \code{dayan_cluster_barrier_t} & 16/8 & arrivals/reserved 位于 0/8。 \\
  \code{dayan_kernel_resource_desc_t} & 32/8 & 16-bit max children/max tasks 为 0/2，32-bit flags 为 4，四类 TCSR count 为 8--11，打包的 Local/Shared count 为 12，atomic cache/URAM size 为 16/24。 \\
  \code{dayan_thread_resources_t} & 8/2 & 四类 TCSR count 为 0--3，Local/Shared count 为 4/5，reserved 为 6。 \\
\end{dayantable}

字段布局属于 ABI；Named Resource 状态只存在于硬件表中。

Resource descriptor 的已知 flag 为
\code{DAYAN_KERNEL_RESOURCE_LOCKED_TASK}。它只允许 max tasks 为 0 或 1，并在 Package
中编码为 \code{NPU_PACKAGE_KERNEL_LOCKED_TASK}；带该位的 Package kernel 只声明
single-Core task mode。Package/header 与 kernel record 当前均为 v5，Runtime/Server
socket 协议为 v10，Runtime options/handshake 为 v5；结构尺寸仍分别保持 128/96 bytes
和 64-byte kernel info。旧版本与新 Server 在握手或 package 校验阶段显式失败。

每个 Hart 的算子 proc ABI 上限为 16 个 slot，root 占一个，因此
\code{DAYAN\_KERNEL\_MAX\_CHILD\_PROCS} 为 15。init、RPC 与 profile proc 不从该池分配。

proc profile 是紧密排列的变长字节流。普通记录为 8 bytes：低 48 bit counter、
8-bit slot 和 8-bit state；\code{SLOT_DEFINE} 记录为 24 bytes，在同一 header 后追加
固定 16-byte proc 名称。名称必须在字段内以 NUL 结尾。状态编码与
\code{THREAD_RUNNING--THREAD_KILL_WAIT_PRODUCE} 的 0--7 一致；Host 将一个 counter
tick 解释为 1 ns。RPC descriptor 的 count/capacity 均以 byte 为单位。

\subsection{目标文件与符号解析}

嵌入的 Device object 必须是 64-bit little-endian RISC-V \code{ET_REL}。加载器支持
当前 dycc 代码模型产生的绝对、PC-relative、CALL/JAL/BRANCH 和 ADD/SUB/SET
relocation；拒绝 TLS、GP-relative、动态链接、构造器 section 和其它 relocation。

\begin{dayancode}
dycc --save-temps=build/dycc -c operator.cpp -o operator.o
llvm-readelf --file-header --relocations \
  build/dycc/operator.cpp.dycc.device.o
llvm-nm --undefined-only build/dycc/operator.cpp.dycc.device.o
\end{dayancode}

未定义符号必须全部存在于 \code{npu-os-api.elf}。\code{dayan-api.symbols} 是精确
白名单；完整固件中的 \code{proc_*}、\code{rpc_*} 和 platform 内部符号不构成算子
ABI。

公开 DAPI 与 libc 分别放入 \code{.text.api.dapi} 和 \code{.text.api.libc}。新增公开
符号时，应同时更新声明、实现、section attribute、白名单、测试和本手册。

\end{document}
